\section{The complex gamma estimate needed for counting}
\label{sec:stirling}

The real factorial Stirling estimate alone is insufficient for the
argument of $\Gamma(1/4+\ii T/2)$. We derive the needed complex estimate
directly. Let $\Log z$ be the branch with argument in
$(-\pi/2,\pi/2)$ on $\Re z>0$, and let $\log\Gamma z$ denote its
analytic logarithm which is real on the positive real axis. It can be
defined by the logarithm of the finite products in
\eqref{eq:gamma-product}; all factors have positive real part after
the initial $z$ factor.

\begin{lemma}[Complex Euler--Maclaurin remainder]
For $\Re z>0$,
\begin{equation}
 \log\Gamma z=(z-\tfrac12)\Log z-z+C+R(z),\qquad C\in\R,
 \label{eq:complex-stirling}
\end{equation}
where
\begin{equation}
 R(z)=\frac1{12z}-\frac12\int_0^\infty
          \frac{B_2(\{x\})}{(z+x)^2}\,dx,
 \qquad |R(z)|<\frac1{4|z|}.
 \label{eq:gamma-remainder}
\end{equation}
The constant is $C=\tfrac12\log(2\pi)$, but its value is not required
in this argument.
\end{lemma}

\begin{proof}
Apply two integrations by parts on each unit interval to
$f(x)=\Log(z+x)$. Since $f'=1/(z+x)$ and $f''=-1/(z+x)^2$, the finite
Euler--Maclaurin identity on $[0,n]$ is
\begin{align*}
 \sum_{j=0}^{n-1}\Log(z+j)
 ={}&\int_0^n\Log(z+x)\,dx
     +\frac{\Log z-\Log(z+n)}2\\
 &+\frac1{12}\left(\frac1{z+n}-\frac1z\right)
   +\frac12\int_0^n\frac{B_2(\{x\})}{(z+x)^2}\,dx.
\end{align*}
Insert this equality into
\[
 \log\Gamma z=
 \lim_{n\to\infty}\left(\log((n-1)!)+z\log n
                 -\sum_{j=0}^{n-1}\Log(z+j)\right).
\]
The elementary antiderivative of $\Log(z+x)$ yields
\eqref{eq:complex-stirling}--\eqref{eq:gamma-remainder}; the
remaining constant is the real limit of
$\log((n-1)!)-(n-1/2)\log n+n$. Existence with an $O(1/n)$ modulus
follows by the same finite formula at $z=1$. No complex limiting
constant has to be selected. The upper-endpoint errors are bounded
uniformly on $|z|\le D$, $\Re z\ge a>0$ by a rational multiple of
$(D+1)^2/n$, using the finite series for $\Log(1+z/n)$ when $n\ge2D$.

Since $|B_2(x)|\le1/6$ on $[0,1]$, it remains to bound a real
integral. Write $z=|z|e^{\ii\theta}$, $|\theta|<\pi/2$. Direct
integration gives, including the limiting case $\theta=0$,
\[
 \int_0^\infty\frac{dx}{|z+x|^2}
 =\frac{|\theta|}{|z|\sin|\theta|}
 \le\frac{\pi}{2|z|}.
\]
Therefore
\[
 |R(z)|\le\frac1{12|z|}+\frac{\pi}{24|z|}
       =\frac{2+\pi}{24|z|}<\frac1{4|z|}.
\]
All divisions are bounded away from zero on the stated half-plane
boxes. For the integral tail beyond $X$ one may use
$1/[12(X+a)]$. On each unit interval, the derivative of its integrand
is bounded by a rational expression in $a^{-1}$ and $|z|$; splitting
at the integer endpoints and taking a fixed rational mesh bounds
quadrature error. Thus the displayed proof has explicit primitive
recursive cutoff and mesh choices. It does not require an unspecified
asymptotic expansion.
\end{proof}

Define the continuous gamma argument
\begin{equation}
 \vartheta(T)=\Im\log\Gamma(\tfrac14+\tfrac{\ii T}2)
               -\frac T2\log\pi,
 \label{eq:theta-definition}
\end{equation}
with the branch just fixed. In particular $\vartheta(0)=0$.
The distinction between this continuous logarithm and the principal
argument of the gamma value matters: reducing modulo $2\pi$ would
destroy a zero-counting formula. The normalization agrees with the
standard Riemann--Siegel theta function; see
\cite{DLMFgamma,DLMFzeros}.

\begin{lemma}[Explicit theta estimate]
For $T\ge1$,
\begin{equation}
 \left|\vartheta(T)-
 \left(\frac T2\log\frac{T}{2\pi}-\frac T2-\frac\pi8\right)\right|
 <\frac1T.
 \label{eq:theta-bound}
\end{equation}
\end{lemma}

\begin{proof}
Put $a=1/4$, $b=T/2$, $r=(a^2+b^2)^{1/2}$ and
$\theta=\arg(a+\ii b)$. The imaginary part of the main expression
in~\eqref{eq:complex-stirling} is
\[
 b\log r-\tfrac14\theta-b.
\]
The replacement of $\log r$ by $\log b$ costs at most
\[
 \frac b2\log(1+a^2/b^2)\le\frac{a^2}{2b}=\frac1{16T}.
\]
The replacement of $\theta$ by $\pi/2$ costs at most
$\tfrac14\arctan(a/b)\le1/(8T)$. Finally
$|R(a+\ii b)|<1/(4b)=1/(2T)$. The total is less than
$11/(16T)<1/T$, and the real constant $C$ contributes no imaginary
part.
\end{proof}
